\documentclass[11pt]{article}
\usepackage[a4paper,margin=27mm]{geometry}
\usepackage{amsmath,amssymb,amsthm}
\usepackage[hidelinks]{hyperref}
\newtheorem{theorem}{Theorem}
\newtheorem{lemma}{Lemma}
\newcommand{\C}{\mathbb C}
\newcommand{\R}{\mathbb R}
\newcommand{\N}{\mathbb N}
\newcommand{\dimH}{\dim_{\mathrm H}}
\title{A counterexample to the dimension clause\protect\\of conjecture 00000007768}
\author{Local formalization artifact}
\date{8 October 2026}
\begin{document}
\maketitle

The original asserts that a complex polynomial outside the power and signed
Chebyshev families has Julia-set fractal dimension strictly greater than one.
We use the standard meanings of affine polynomial conjugacy and Hausdorff
dimension. The following exact example refutes this necessary clause of the
conjecture. It does not require deciding the separate algebraic-arc,
transcendence or cusp-parameter assertions.

For a polynomial $p$ of degree at least two, write
\[
 K(p)=\{z\in\C:\text{the complete forward orbit of }z\text{ is bounded}\},
 \qquad J(p)=\partial K(p).
\]
This is the usual filled Julia set and polynomial Julia set; see Milnor,
\emph{Dynamics in One Complex Variable}, \S17, Lemma 17.1. All boundedness below
is over every iterate, not over a finite truncation.

\begin{theorem}
Let $f(z)=z^2-6$. Its coefficients are algebraic, its degree is two, and it is
not affinely conjugate to $z^d$, $T_d$, or $-T_d$ for any $d\in\N$. Its Julia
set is nonempty and satisfies $J(f)\subset\R\subset\C$, hence
$\dimH J(f)\leq 1$. Thus the original universal assertion
``non-exceptional implies dimension $>1$'' is false.
\end{theorem}

\begin{lemma}[Escape radius, with the full orbit]
An orbit of $f$ is bounded if and only if all its terms have modulus at most
three. Consequently $K(f)$ is closed.
\end{lemma}
\begin{proof}
The reverse triangle inequality gives
\[
 |f(z)|\geq |z|^2-6.
\]
Suppose an orbit term has modulus $r_0>3$, and put $\delta=r_0-3>0$.
For every $r\geq r_0$,
\[
 r^2-6-r=(r-3)(r+2)\geq 5\delta\geq\delta.
\]
Induction therefore bounds the subsequent moduli below by $r_0+n\delta$.
The Archimedean property makes this unbounded, a contradiction. The argument
applies to every tail of a bounded orbit, so every term has modulus at most
three. The converse is immediate. Moreover
\[
 K(f)=\bigcap_{n\in\N}\{z:|f^{\circ n}(z)|\leq 3\}.
\]
Each set in the intersection is closed by continuity, proving the last claim.
\end{proof}

\begin{lemma}[Reality and nonemptiness]
$K(f)\subset\R$ and $3\in J(f)$.
\end{lemma}
\begin{proof}
Let $z_n=x_n+iy_n$ be a bounded forward orbit. The escape lemma gives
$|z_{n+1}|\leq 3$, so
\[
 x_n^2-y_n^2-6=\operatorname{Re} z_{n+1}\geq -3,
 \qquad x_n^2-y_n^2\geq 3.
\]
In particular $|x_n|\geq 1$. Since $y_{n+1}=2x_ny_n$,
\[
 |y_n|\geq 2^n|y_0|.
\]
If $y_0\ne0$, this contradicts $|y_n|\leq |z_n|\leq3$. Hence every point of
$K(f)$ is real. Because the real axis is closed,
\[
 J(f)=\partial K(f)\subset\overline{K(f)}\subset\R.
\]
The point $3$ is fixed, so belongs to $K(f)$. Every neighborhood of $3$ contains
a nonreal point, and such a point is outside $K(f)$ by the reality argument.
Therefore $3$ belongs to the boundary $J(f)$, proving nonemptiness.
\end{proof}

\begin{lemma}[Every exceptional degree and both Chebyshev signs]
The polynomial $f$ is not affinely conjugate to $z^d$ or $\pm T_d$ for any
$d\in\N$.
\end{lemma}
\begin{proof}
Write a putative affine conjugacy as $h(z)=az+b$, with $a\ne0$, satisfying
$f\circ h=h\circ q$. This map is bijective, with inverse $(z-b)/a$.
Polynomial degrees in this identity give
\[
 2=\deg q=d,
\]
using the actual degree of $T_d$, and the same degree for its negative. Thus
all degrees reduce to $d=2$.

Now $q(1)=q(-1)$ for each of $q(z)=z^2$, $T_2(z)=2z^2-1$, and $-T_2(z)$.
Evaluating $f(h(z))=h(q(z))$ at $1$ and $-1$ and subtracting gives $4ab=0$.
As $a\ne0$, we have $b=0$.

If $q(z)=z^2$, evaluation at zero gives $-6=0$, impossible. If
$q(z)=2z^2-1$, the zero evaluation forces $a=6$; if
$q(z)=-(2z^2-1)$, it forces $a=-6$. In either case, evaluation at one then
gives $36-6=6$, again impossible. This uses the actual signed Chebyshev
polynomials, not a substituted centered-parameter hypothesis.
\end{proof}

\begin{proof}[Completion of the theorem]
The integer coefficients of $z^2-6$ are algebraic over $\mathbb Q$.
The real-axis inclusion above and monotonicity of Hausdorff dimension give
\[
 \dimH J(f)\leq\dimH\R=1,
\]
where the real axis is identified with its isometric image in the complex
plane. Together with the complete exceptional-family exclusion and
nonemptiness, this is the claimed counterexample. A conjunction containing
the asserted strict lower bound therefore cannot hold for all the original
admissible polynomials.
\end{proof}

\paragraph{Source and formalization correspondence.}
Both the English and Chinese source copies assert the dimension clause used
here. The Lean objects use polynomials over $\C$, arbitrary natural-number
iterates, bounded-orbit filled Julia sets, their topological boundaries, actual
Mathlib Hausdorff dimension, actual Mathlib Chebyshev polynomials, and
invertible affine conjugacies. The named dimension convention and conjugacy
convention are explicit semantic conventions and remain subject to independent
source review. Compiler checks of an author package do not constitute manager
acceptance or an external submission.

\begin{thebibliography}{1}
\bibitem{milnor}
John Milnor, \emph{Dynamics in One Complex Variable: Introductory Lectures},
Stony Brook IMS Preprint 1990/5, revised September 1991,
\S17, Lemma 17.1 (PDF page 98).
\url{https://legacy-www.math.harvard.edu/archive/118r_spring_05/docs/milnor.pdf}
\end{thebibliography}
\end{document}
